\ifdefined\PRINT
%  % 扉画像ページ(印刷版)
%  % 印刷版PDFには表紙画像を含めず、本文とは別に入稿するのが基本
%  % 表紙をめくったら1ページ目に扉としてモノクロ版の画像を入れるような場合には
%  % コメントアウトされている部分を生かす
%  \thispagestyle{empty} % emptyでも隠しノンブルは入る
%  % トンボ(1in)と塗り足し幅(3mm)のぶんだけ位置がズレるので補正
%  % 1inという値はjlreq-trimmarksのデフォルトが2in大きい用紙を使うことに由来
%  % デフォルトから変更していたり、別の方法でトンボを作る場合には要修正
%  \setlength{\wpXoffset}{22.4mm}  % 1in-3mm
%  \setlength{\wpYoffset}{-28.4mm} % -1in-3mm
%  % 画像は版面いっぱいではなく、塗り足し幅(3mm)の分拡大する
%  % (A5幅148mm + 塗り足し幅3mm x 2) / 148mm = 1.041; B5なら(182+6)/182=1.033
%  % 画像ファイルは縦横比1:1.414ではなく、154:216 (B5なら188:263)の比率で
%  % グレースケール(またはCMYK)の高解像度のものを指定する
%  % A5グレースケール600dpiドットバイドットなら3638x5102ピクセル
%  \ThisCenterWallPaper{1.041}{cover/print.png}
%  \null\clearpage
\else
  % 表紙(電子版)
  % \ThisCenterWallPaper の最初の引数は版面の横幅に対する画像幅の比率
  % 1なら版面にちょうどぴったり収まるよう画像サイズを調整する
  % 縦横比は保持されるので、1:1.414 の比率の画像ファイルを指定すると
  % 縦もぴったり版面に収まる
  %\setcounter{page}{0} % 表紙画像をページ数に数えない場合
  \thispagestyle{empty}
  \ThisCenterWallPaper{1}{cover/cover.png}
  \null% この\nullがないと表示されない
  \clearpage
\fi

% 序文
\chapter{序}

なんかかっこいいことを書く。

